\documentclass[a4paper,12pt]{article}

\usepackage[T2A]{fontenc}
\usepackage[utf8]{inputenc}
\usepackage[russian]{babel}

\usepackage{amsmath, amssymb, amsthm, mathtools}
\usepackage{helvet}
\renewcommand{\familydefault}{\sfdefault}
\usepackage{icomma}
\usepackage[left=18mm,right=18mm,top=18mm,bottom=20mm]{geometry}
\usepackage{microtype}
\usepackage{tikz}
\usepackage{tkz-euclide}
\usepackage{pgfplots}
\pgfplotsset{compat=1.18}
\usepackage{tabularx, booktabs, longtable}
\usepackage{enumitem}
\usepackage{hyperref}
\usepackage{parskip}
\usepackage{fancyhdr}
\usepackage{setspace}
\usepackage{xcolor}

\definecolor{accent}{HTML}{008080}
\definecolor{darkaccent}{HTML}{004D4D}
\definecolor{textgray}{HTML}{333333}
\definecolor{softgray}{HTML}{F3F5F5}
\definecolor{lightaccent}{HTML}{E6F2F2}

\hypersetup{
    colorlinks=true,
    linkcolor=darkaccent,
    urlcolor=darkaccent
}

\onehalfspacing
\setlength{\parindent}{0pt}
\setlength{\parskip}{6pt}

\pagestyle{fancy}
\fancyhf{}
\fancyhead[L]{\small Чек-лист}
\fancyhead[R]{\small Системы и совокупности условий}
\fancyfoot[C]{\small \thepage}
\renewcommand{\headrulewidth}{0.4pt}
\renewcommand{\headrule}{\hbox to\headwidth{\color{accent}\leaders\hrule height \headrulewidth\hfill}}

\setlist[itemize]{leftmargin=1.2cm,itemsep=2pt,topsep=4pt}
\setlist[enumerate]{leftmargin=1.2cm,itemsep=4pt,topsep=4pt}

\newcommand{\important}[1]{\textcolor{darkaccent}{\textbf{#1}}}
\newcommand{\ruleline}{\vspace{2mm}{\color{accent}\hrule height 0.5pt}\vspace{2mm}}
\newcommand{\blank}{\underline{\hspace{3cm}}}

\newcommand{\intervalaxis}{
\begin{center}
\begin{tikzpicture}[scale=0.9]
    \draw[->, color=textgray] (-3.2,0) -- (4.2,0);
    \foreach \x/\lab in {-2/$-2$,0/$0$,3/$3$,4/$8$}
        \draw[color=textgray] (\x,0.08) -- (\x,-0.08) node[below] {\small \lab};
    \draw[line width=1.6pt, color=accent] (-2,0.23) -- (4,0.23);
    \filldraw[fill=white, draw=accent, line width=1pt] (-2,0.23) circle (2pt);
    \filldraw[fill=accent, draw=accent] (4,0.23) circle (2pt);
    \draw[line width=1.6pt, color=darkaccent] (3,0.48) -- (4,0.48);
    \filldraw[fill=white, draw=darkaccent, line width=1pt] (3,0.48) circle (2pt);
    \filldraw[fill=darkaccent, draw=darkaccent] (4,0.48) circle (2pt);
    \node[above, color=accent] at (1,0.25) {\small $-2\le x\le 8$};
    \node[above, color=darkaccent] at (3.5,0.55) {\small $x>3$};
\end{tikzpicture}
\end{center}
}

\begin{document}

\begin{titlepage}
\thispagestyle{empty}
\begin{center}
\vspace*{2.2cm}

{\Large \textcolor{accent}{\textbf{Чек-лист}}}

\vspace{8mm}

{\huge \textbf{Когда появляется совокупность условий}}

\vspace{4mm}

{\Large \textbf{Системы, совокупности и логика выбора}}

\vspace{18mm}

{\large ЕГЭ}

\vspace{22mm}

{\large Лёвин Артём Александрович эксклюзивно для Николь Саркисянц}

\vspace{8mm}

{\large \today}

\vfill

\begin{tikzpicture}[remember picture, overlay]
    \draw[line width=1.1pt, color=accent]
        ([xshift=2.3cm,yshift=2.2cm]current page.south west)
        .. controls ([xshift=6.0cm,yshift=3.3cm]current page.south west)
        and ([xshift=-6.0cm,yshift=1.3cm]current page.south east)
        .. ([xshift=-2.3cm,yshift=2.4cm]current page.south east);
    \draw[line width=0.5pt, color=darkaccent]
        ([xshift=3.0cm,yshift=1.55cm]current page.south west)
        .. controls ([xshift=7.4cm,yshift=2.25cm]current page.south west)
        and ([xshift=-7.0cm,yshift=2.7cm]current page.south east)
        .. ([xshift=-3.0cm,yshift=1.65cm]current page.south east);
\end{tikzpicture}

\vspace*{1.2cm}
\end{center}
\end{titlepage}

\section*{1. Главная идея}

\important{Система} условий означает, что должны выполняться \important{все} условия одновременно. В записи это часто изображают фигурной скобкой:
\[
\begin{cases}
A,\\
B.
\end{cases}
\]
Математически это означает пересечение множеств:
\[
A\cap B.
\]

\important{Совокупность} условий означает, что достаточно выполнения \important{хотя бы одного} условия. В школьной записи это часто изображают квадратной скобкой:
\[
\left[
\begin{array}{l}
A,\\
B.
\end{array}
\right.
\]
Математически это означает объединение множеств:
\[
A\cup B.
\]

\section*{2. Короткое правило выбора}

\begin{center}
\renewcommand{\arraystretch}{1.35}
\begin{tabularx}{\linewidth}{>{\bfseries}p{0.28\linewidth}X}
\toprule
Ситуация & Что использовать \\
\midrule
ОДЗ: логарифм, корень чётной степени, знаменатель дроби & Система: ограничения должны выполняться одновременно с основным условием. \\
Модуль раскрывается по случаям & Совокупность или разбиение на случаи: подходит один из вариантов. \\
Произведение равно нулю & Совокупность: хотя бы один множитель равен нулю. \\
Дробь равна нулю & Система внутри случая: числитель равен нулю, знаменатель не равен нулю. \\
Квадратное уравнение разложено на множители & Совокупность: каждый множитель может дать свой корень. \\
Уравнение с чётной степенью & Часто возникает совокупность: например, $u^2=a^2 \Rightarrow u=a$ или $u=-a$. \\
Неравенство с несколькими ограничениями & Обычно система: все ограничения нужно пересечь с решением неравенства. \\
\bottomrule
\end{tabularx}
\end{center}

\section*{3. Система: когда нужны все условия}

Система появляется тогда, когда решение должно удовлетворять сразу нескольким требованиям.

\subsection*{Пример с логарифмом}

Рассмотрим условия:
\[
\begin{cases}
|x-3|\le 5,\\
x-3>0.
\end{cases}
\]

Первое условие:
\[
|x-3|\le 5
\quad\Longleftrightarrow\quad
-5\le x-3\le 5
\quad\Longleftrightarrow\quad
-2\le x\le 8.
\]

Второе условие:
\[
x-3>0
\quad\Longleftrightarrow\quad
x>3.
\]

Так как это система, берём пересечение:
\[
[-2;8]\cap(3;+\infty)=(3;8].
\]

\intervalaxis

\[
\boxed{x\in(3;8]}
\]

\important{Вывод.} Если одно условие пришло из ОДЗ, а другое из самого уравнения или неравенства, то почти всегда нужно пересекать ответы.

\section*{4. Совокупность: когда достаточно одного варианта}

Совокупность появляется тогда, когда условие может выполниться разными независимыми способами.

\subsection*{Случай 1. Модуль равен числу}

\[
|x-3|=5.
\]

Модуль равен $5$, если выражение под модулем равно $5$ или $-5$:
\[
\left[
\begin{array}{l}
x-3=5,\\
x-3=-5.
\end{array}
\right.
\]

Отсюда:
\[
\left[
\begin{array}{l}
x=8,\\
x=-2.
\end{array}
\right.
\]

Ответ:
\[
\boxed{x\in\{-2;8\}}.
\]

\important{Почему это совокупность?} Число подходит, если выполняется хотя бы одно из двух равенств.

\subsection*{Случай 2. Произведение равно нулю}

\[
(x-4)(x+7)=0.
\]

Произведение равно нулю, если хотя бы один множитель равен нулю:
\[
\left[
\begin{array}{l}
x-4=0,\\
x+7=0.
\end{array}
\right.
\]

\[
\left[
\begin{array}{l}
x=4,\\
x=-7.
\end{array}
\right.
\]

Ответ:
\[
\boxed{x\in\{-7;4\}}.
\]

\important{Ловушка.} Нельзя писать систему:
\[
\begin{cases}
x-4=0,\\
x+7=0.
\end{cases}
\]
Такая система потребовала бы, чтобы $x$ одновременно равнялся $4$ и $-7$, что невозможно.

\subsection*{Случай 3. Квадратное уравнение после разложения}

\[
x^2-5x+6=0.
\]

Разложим:
\[
x^2-5x+6=(x-2)(x-3).
\]

Получаем:
\[
(x-2)(x-3)=0.
\]

Значит:
\[
\left[
\begin{array}{l}
x-2=0,\\
x-3=0.
\end{array}
\right.
\]

Ответ:
\[
\boxed{x\in\{2;3\}}.
\]

\subsection*{Случай 4. Чётная степень}

\[
(x-1)^2=16.
\]

Так как $16=4^2$, получаем:
\[
x-1=4
\quad\text{или}\quad
x-1=-4.
\]

Записываем совокупность:
\[
\left[
\begin{array}{l}
x-1=4,\\
x-1=-4.
\end{array}
\right.
\]

\[
\left[
\begin{array}{l}
x=5,\\
x=-3.
\end{array}
\right.
\]

Ответ:
\[
\boxed{x\in\{-3;5\}}.
\]

\section*{5. Смешанные ситуации: совокупность внутри системы}

В задачах ЕГЭ часто встречаются смешанные случаи: сначала появляется совокупность, но каждый её вариант должен ещё пройти проверку ОДЗ.

\subsection*{Пример}

\[
\frac{(x-2)(x+5)}{x-2}=0.
\]

Сначала записываем ОДЗ:
\[
x-2\ne0
\quad\Longleftrightarrow\quad
x\ne2.
\]

Дробь равна нулю, когда числитель равен нулю, а знаменатель не равен нулю:
\[
\begin{cases}
(x-2)(x+5)=0,\\
x-2\ne0.
\end{cases}
\]

Внутри числителя появляется совокупность:
\[
\left[
\begin{array}{l}
x-2=0,\\
x+5=0.
\end{array}
\right.
\]

Кандидаты:
\[
x=2,\qquad x=-5.
\]

Проверяем ОДЗ:
\[
x=2 \quad \text{не подходит, так как знаменатель равен нулю;}
\]
\[
x=-5 \quad \text{подходит.}
\]

Ответ:
\[
\boxed{x=-5}.
\]

\important{Ключевой вывод.} Совокупность даёт кандидатов, а система с ОДЗ отсекает лишние значения.

\section*{6. Модуль: равенство, неравенство и случаи}

\subsection*{Модуль равен положительному числу}

\[
|A|=b,\qquad b>0.
\]

Тогда:
\[
\left[
\begin{array}{l}
A=b,\\
A=-b.
\end{array}
\right.
\]

Это совокупность, потому что подходят два независимых варианта.

\subsection*{Модуль меньше или равен числу}

\[
|A|\le b,\qquad b\ge0.
\]

Тогда:
\[
-b\le A\le b.
\]

Это двойное неравенство, то есть фактически система:
\[
\begin{cases}
A\ge -b,\\
A\le b.
\end{cases}
\]

\subsection*{Модуль больше или равен числу}

\[
|A|\ge b,\qquad b\ge0.
\]

Тогда:
\[
\left[
\begin{array}{l}
A\ge b,\\
A\le -b.
\end{array}
\right.
\]

Это совокупность, потому что выражение может быть либо достаточно большим положительным, либо достаточно большим отрицательным.

\begin{center}
\renewcommand{\arraystretch}{1.35}
\begin{tabularx}{\linewidth}{p{0.25\linewidth}p{0.35\linewidth}X}
\toprule
Вид & Переход & Логика \\
\midrule
$|A|=b$, $b>0$ & $A=b$ или $A=-b$ & Совокупность \\
$|A|\le b$, $b\ge0$ & $-b\le A\le b$ & Пересечение \\
$|A|\ge b$, $b\ge0$ & $A\ge b$ или $A\le -b$ & Совокупность \\
\bottomrule
\end{tabularx}
\end{center}

\section*{7. Типовые ошибки}

\begin{enumerate}
    \item \important{Путать систему и совокупность.}

    Неверно:
    \[
    (x-1)(x-4)=0
    \quad\Rightarrow\quad
    \begin{cases}
    x-1=0,\\
    x-4=0.
    \end{cases}
    \]
    Верно:
    \[
    (x-1)(x-4)=0
    \quad\Rightarrow\quad
    \left[
    \begin{array}{l}
    x-1=0,\\
    x-4=0.
    \end{array}
    \right.
    \]

    \item \important{Забывать ОДЗ.}

    В выражении
    \[
    \log_2(x-3)
    \]
    обязательно:
    \[
    x-3>0.
    \]

    \item \important{Сокращать дробь без учёта запрещённых значений.}

    Если
    \[
    \frac{(x-2)(x+1)}{x-2}=0,
    \]
    то значение $x=2$ запрещено, даже если множитель сократился.

    \item \important{Неверно раскрывать модуль.}

    Для равенства $|A|=b$ нужна совокупность, а для $|A|\le b$ нужен промежуток от $-b$ до $b$.

    \item \important{Не проверять кандидаты.}

    После совокупности всегда полезно проверить найденные значения в исходном выражении, особенно если были дроби, корни, логарифмы или возведение в квадрат.
\end{enumerate}

\section*{8. Мини-алгоритм решения}

\begin{enumerate}
    \item Найдите ОДЗ, если есть логарифмы, знаменатели, корни чётной степени.
    \item Решите основное уравнение или неравенство.
    \item Определите, где нужна система, а где совокупность.
    \item Если появились несколько ветвей, решите каждую ветвь отдельно.
    \item Объедините ответы из совокупности.
    \item Пересеките полученный результат с ОДЗ.
    \item Проверьте спорные значения в исходном выражении.
\end{enumerate}

\newpage

\section*{Тренировочный блок}

\textbf{Инструкция.} Решите упражнения. В каждом задании отдельно отметьте, где нужна система, а где совокупность. Ответы запишите в виде числа, множества или промежутка.

\subsection*{Средний уровень}

\begin{enumerate}
    \item Решите уравнение:
    \[
    |x-4|=7.
    \]

    \item Решите уравнение:
    \[
    (x-5)(x+2)=0.
    \]

    \item Решите систему условий:
    \[
    \begin{cases}
    |x-1|\le 4,\\
    x-1>0.
    \end{cases}
    \]
\end{enumerate}

\subsection*{Выше среднего уровня}

\begin{enumerate}[resume]
    \item Решите уравнение:
    \[
    (x-3)^2=25.
    \]

    \item Решите неравенство:
    \[
    |x+2|\ge 6.
    \]

    \item Решите уравнение:
    \[
    \frac{(x-4)(x+1)}{x-4}=0.
    \]

    \item Решите уравнение:
    \[
    \sqrt{x+6}=x.
    \]

    \item Решите неравенство:
    \[
    \log_2(x-1)\le 3.
    \]

    \item Решите уравнение:
    \[
    \frac{x^2-9}{x-3}=0.
    \]
\end{enumerate}

\subsection*{Сложный уровень}

\begin{enumerate}[resume]
    \item Решите неравенство:
    \[
    \log_3|x-2|\le 2.
    \]
\end{enumerate}

\newpage

\section*{Ответы для самопроверки}

\begin{center}
\renewcommand{\arraystretch}{1.45}
\begin{tabularx}{\linewidth}{>{\bfseries}p{0.14\linewidth}X}
\toprule
№ & Ответ \\
\midrule
1 & $x\in\{-3;11\}$ \\
2 & $x\in\{-2;5\}$ \\
3 & $x\in(1;5]$ \\
4 & $x\in\{-2;8\}$ \\
5 & $x\in(-\infty;-8]\cup[4;+\infty)$ \\
6 & $x=-1$ \\
7 & $x=3$ \\
8 & $x\in(1;9]$ \\
9 & $x=-3$ \\
10 & $x\in[-7;2)\cup(2;11]$ \\
\bottomrule
\end{tabularx}
\end{center}

\end{document}